\documentclass[10pt,a4paper]{amsart}
\usepackage[margin=19mm]{geometry}
\usepackage{amsmath,amssymb,mathtools}
\usepackage[scr=rsfs,cal=euler]{mathalfa}
\usepackage{xfrac}
\usepackage[unicode,hidelinks,bookmarks=false]{hyperref}
\newcommand{\ie}{i.e.\ }
\newcommand{\cf}{cf.\ }
\newcommand{\resp}{respectively\ }
\newcommand{\Subsection}{\subsection}
\providecommand{\nobreakdash}{\nobreak\hbox{-}}
\setlength{\emergencystretch}{2em}
\multlinegap0pt
\usepackage{amssymb}
\usepackage{mathtools}
\usepackage[shortlabels]{enumitem}
\newtheorem{theorem}{Theorem}
\usepackage{cleveref}
\crefname{theorem}{Theorem}{Theorems}
\newcommand{\D}{\mathcal{D}}
\newcommand{\Lup}{\operatorname{Lip}^{+}_{1}(\R)}
\newcommand{\R}{\mathbb{R}}
\newcommand{\Xp}{\mathcal{X}^{+}}
\DeclareMathOperator{\lipone}{Lip_1}
\newcommand{\lipplus}{\operatorname{Lip}^{+}_{1}}
\title{Pyramidal Compactification of Asymmetric Metric Measure Spaces via Adjoint Transport}
\author{Shigeaki Yokota}
\date{Source: arXiv:2608.01145v1 (CC BY 4.0)}
\begin{document}
\maketitle

\noindent\emph{Source and licence.} Pyramidal Compactification of Asymmetric Metric Measure Spaces via Adjoint Transport. Version arXiv:2608.01145v1 (CC BY 4.0), distributed by arXiv under the Creative Commons Attribution 4.0 International licence, http://creativecommons.org/licenses/by/4.0/, as recorded in the arXiv licence metadata for this exact version and shown on the abstract page https://arxiv.org/abs/2608.01145v1. Full licence notice: \texttt{licenses/arxiv-2608.01145-CC-BY-4.0.txt}.\par\smallskip

\noindent\textit{Editorial scope.} Counted unit: the article's theorem thm:adjunction-direction (source0.tex lines 377--385). The article's complete original proof of it is delivered with it (source0.tex lines 387--405, one \texttt{proof} environment of 19 lines). No mathematics is written, repaired or re-derived: the statement and the proof are the article's own lines, in the article's own notation, and no retained line is substituted or re-flowed.  Retained uncounted as the source-local context the proof needs: lines 263--266.  Nothing was omitted from any retained span, so the package carries no omission record.  No retained line contains a citation, so the wrapper carries no bibliography and no bibliography entry.  No retained line contains a figure environment, an image-inclusion command or drawing code, so no image is delivered and the package declares no asset dependency.  Editorial formatting only, in addition: an AMS \texttt{amsart} preamble that restates the article's own declarations and macros used by the retained lines, namely the environments \texttt{theorem} on separate counters, together with the standard packages those lines need.  The retained environments print in this document as Theorem 1 in order of appearance, whereas the article prints the same items with the numbering of its own full document.  The complete native source of the article as distributed in the arXiv e-print for this exact version is bundled unchanged as \texttt{sources/206599/source0.tex}.
\begin{align}
d_X(x,x')&=\sup_{f\in\lipone(X)}|f(x)-f(x')|,\label{eq:symmetric-recovery}\\
d_Y(y,y')&=\sup_{f\in\lipplus(Y)}\{f(y')-f(y)\}.\label{eq:directed-recovery}
\end{align}

\begin{theorem}[Representation--reconstruction adjunction for qm-spaces]
\label{thm:adjunction-direction}
There is a natural adjunction $\operatorname{Rep}^{+}\dashv\operatorname{Rec}^{+}$. More precisely, there is a natural bijection that leaves the underlying maps unchanged,
\begin{equation}
\label{eq:hom-asymmetric}
\operatorname{Hom}_{\Lup\circ\D}(\operatorname{Rep}^{+}(Y),Y')\simeq\operatorname{Hom}_{\Xp}(Y,\operatorname{Rec}^{+}(Y')).
\end{equation}
The unit of this adjunction is an isomorphism, and $\operatorname{Rep}^{+}(-)$ is fully faithful.
\end{theorem}

\begin{proof}
\emph{Claim.} The assignment $\operatorname{Rec}^{+}(-)$ is a functor from $\Lup\circ\D$ to $\Xp$.

The family $\overline{F_X}$ is $\Lup$-closed, and $\Lup$ contains every constant function. Fix $f_0\in\overline{F_X}$. Every constant function belongs to $\overline{F_X}$ as the composition of a constant map with $f_0$. Therefore, $d_X^+$ is nonnegative and $d_X^+(x,x)=0$. For every $f\in\overline{F_X}$, the decomposition
\[f(x'')-f(x)=\{f(x'')-f(x')\}+\{f(x')-f(x)\}\]
and taking suprema give the triangle inequality for $d_X^+$. We also have
\[\max\{d_X^+(x,x'),d_X^+(x',x)\}=\sup_{f\in\overline{F_X}}|f(x)-f(x')|=d_{F_X}(x,x').\]
Thus, the symmetrization agrees with the metric of the gd-set, and the completeness, separability, and measure conditions also hold.

Let $u\colon X\to Y$ be a domination of gd-sets. Taking the pointwise closure of $F_Y\circ u\subset\overline{F_X}$ gives $\overline{F_Y}\circ u\subset\overline{F_X}$. Taking the supremum of the increments therefore yields $d_Y^+(u(x),u(x'))\leq d_X^+(x,x')$. Thus, $\operatorname{Rec}^{+}(-)$ maps morphisms to morphisms.
This proves the claim.

Let $u\colon Y\to Y'$ be a measure-preserving map. Since $\lipplus(Y)$ is closed under pointwise convergence, $u$ is a morphism of gd-sets from $\operatorname{Rep}^{+}(Y)$ to $Y'$ if and only if
\[F_{Y'}\circ u\subset\lipplus(Y).\]
If this inclusion holds, then $\overline{F_{Y'}}\circ u\subset\lipplus(Y)$ because $\lipplus(Y)$ is closed under pointwise convergence. Taking the supremum of the increments, we obtain $d_{Y'}^+(u(y),u(y'))\leq d_Y(y,y')$. Conversely, if $u$ is $1$-Lipschitz, then every $f\in F_{Y'}$ satisfies
\[f(u(y'))-f(u(y))\leq d_{Y'}^+(u(y),u(y'))\leq d_Y(y,y').\]
This proves \Cref{eq:hom-asymmetric}. The bijection leaves the underlying maps unchanged and is therefore natural. The unit $Y\to\operatorname{Rec}^{+}(\operatorname{Rep}^{+}(Y))$ is the identity map and is an isomorphism by \Cref{eq:directed-recovery}. Therefore, $\operatorname{Rep}^{+}$ is fully faithful.
This completes the proof.
\end{proof}

\end{document}
